% ====================================================================================
\chapter{信息与编码：最大熵原理}
% ====================================================================================

\section*{从猜数字游戏说起}

假设我心里想了一个1到8之间的整数，你要用“是/否”问题来猜出它。

\textbf{策略A}（逐个问）：
\begin{quote}
“是1吗？” → 否 → “是2吗？” → 否 → $\cdots$
\end{quote}
最坏情况需要7次。

\textbf{策略B}（二分法）：
\begin{quote}
“大于4吗？” → 是 → “大于6吗？” → 否 → “是5吗？” → 否 → “是6！”
\end{quote}
无论如何只需要3次。

为什么是3次？因为 $\log_2 8 = 3$。每个“是/否”问题最多能把可能性减半，8种可能性需要 $\log_2 8 = 3$ 次二分。

\begin{insightbox}
这个简单的游戏揭示了信息论的核心思想：\textbf{信息量 = 消除不确定性所需的最少问题数}。

Shannon把这个“问题数”叫做\textbf{比特（bit）}——1 bit就是一个“是/否”问题能传递的信息量。
\end{insightbox}

\storynote{Claude Shannon}{1916--2001。他是个有趣的人。他喜欢骑独轮车，会在贝尔实验室的走廊里边骑边抛球杂耍。他发明了“bit”这个词，建立了整个信息论学科。业余时间，他炒股赚了很多钱，还造了各种稀奇古怪的机器。}

\vspace{0.5cm}
\noindent\textit{为什么压缩文件有极限？为什么分类网络的最后一层非用 softmax 不可？为什么“我只知道平均值”这一句话，就足以推出整个分布？}

\begin{quote}
\textbf{本章目标}：你会发现，“信息”这个日常概念可以被精确量化，而计算最大熵分布的方法，正是我们在第1章学过的拉格朗日乘子法！逻辑回归、softmax函数、交叉熵损失，都可以从最大熵原理推导出来。
\end{quote}

\begin{tcolorbox}[colback=green!5, colframe=green!60!black, title=学完这章，你将能够...]
\begin{itemize}
    \item 计算离散分布和连续分布的熵
    \item 用拉格朗日乘子法推导最大熵分布（如指数分布、正态分布）
    \item 解释为什么机器学习用交叉熵作为损失函数
    \item 从最大熵原理推导softmax函数和逻辑回归
    \item 理解KL散度、交叉熵与最大似然之间的关系
\end{itemize}
\end{tcolorbox}

% ------------------------------------------------------------------------------------
\section{信息量的直觉}

\subsection{什么消息更有“信息量”？}

要把“信息”变成一个可以计算的量，先得弄清楚它到底和什么有关。比较两条消息：消息A说“明天太阳从东边升起”，消息B说“明天北京下雪”（假设现在是7月）。哪条消息的信息量更大？直觉告诉我们是消息B，因为消息A是必然发生的事（概率 $\approx 1$），听了等于没听；而消息B是罕见事件（概率很小），一旦成真就大大改变了我们对明天的预期。于是信息量应当是概率的减函数——概率越小，消息越“值钱”。

\begin{keyformula}[title=单个事件的信息量]
事件 $x$ 的信息量（自信息）定义为：
\begin{equation}
I(x) = -\log_2 p(x) \quad \text{（单位：比特）}
\end{equation}
\end{keyformula}
\clarify{这里$p(x)$是事件$x$发生的概率。概率是0到1之间的数，表示事件发生的可能性。比如“抛硬币正面朝上”的概率是0.5，“太阳从东边升起”的概率是1。}

用几个例子验证这个定义是否符合直觉。必然事件（$p=1$）的信息量是 $I = -\log_2 1 = 0$ 比特，没有新信息；抛硬币得到正面（$p=0.5$）的信息量是 $I = -\log_2 0.5 = 1$ 比特；掷骰子得到 6（$p=1/6$）的信息量是 $I = -\log_2(1/6) \approx 2.58$ 比特；一个概率为 $0.001$ 的罕见事件则携带 $I = -\log_2 0.001 \approx 10$ 比特。概率每减半，信息量恰好增加一比特。

\begin{physicalmeaning}
为什么偏偏是 $-\log p$，而不是别的减函数？因为除了“概率越小，信息量越大”（$-\log$ 是单调递减函数）之外，我们对信息量还有两条要求。其一，独立事件的信息量应当可加：同时得知两件互不相关的事，得到的信息应等于分别得知它们的信息之和，而独立事件的概率是相乘的，$I(x \cap y) = -\log(p_x p_y) = -\log p_x - \log p_y = I(x) + I(y)$，对数恰好把乘法变成加法。其二，确定事件的信息量应为零，而 $-\log 1 = 0$。这三条性质唯一地确定了 $-\log p$ 的形式！
\end{physicalmeaning}

\subsection{回到猜数字游戏}

现在我们可以理解为什么猜1-8的数字需要3比特：

每个数字的概率是 $p = 1/8$，信息量是：
\[
I = -\log_2 \frac{1}{8} = \log_2 8 = 3 \text{ 比特}
\]

\textbf{推广}：猜 $n$ 个等可能事件中的一个，需要 $\log_2 n$ 比特。

% ------------------------------------------------------------------------------------
\section{熵：平均信息量}

\subsection{不均匀概率的情况}

猜数字游戏里每个答案等可能，所以每条消息的信息量都一样。现实中的消息源却往往不均匀：假设天气预报说晴天概率 70\%、雨天概率 30\%。如果告诉你“明天是晴天”，信息量是 $-\log_2 0.7 \approx 0.51$ 比特；如果告诉你“明天下雨”，信息量是 $-\log_2 0.3 \approx 1.74$ 比特。两种消息的信息量不同，那么\textbf{平均来看}，每天的天气消息能带来多少信息？自然的做法是按概率加权：

\begin{align}
\text{平均信息量} &= 0.7 \times (-\log_2 0.7) + 0.3 \times (-\log_2 0.3) \\
&= 0.7 \times 0.51 + 0.3 \times 1.74 \\
&\approx 0.88 \text{ 比特}
\end{align}

这个“平均信息量”就是\textbf{熵}。香农 1948 年在《通信的数学理论》里引入这个量时，据说是冯·诺伊曼建议他把它叫作“熵”，理由之一是“没人真正知道熵是什么，所以辩论时你总占上风”。名字的来历下一节会讲清楚——它确实和物理学里的熵是同一个东西。

\begin{keyformula}[title=熵（Shannon Entropy）]
\begin{equation}
H(p) = -\sum_{i=1}^{k} p_i \log_2 p_i = \mathbb{E}[-\log_2 p(X)]
\end{equation}

熵 = 信息量的期望 = 平均每次观测带来的“惊讶程度”
\end{keyformula}
\clarify{$\mathbb{E}[\cdot]$表示“期望”，就是加权平均。如果你没学过概率论，只需记住：期望 = $\sum_i p_i \times (\text{第}i\text{个值})$，即每个值乘以它的概率再求和。}

\subsection{熵的计算例子}

\begin{example}[抛硬币]
公平硬币：$p(\text{正}) = p(\text{反}) = 0.5$
\[
H = -0.5 \log_2 0.5 - 0.5 \log_2 0.5 = 0.5 + 0.5 = 1 \text{ 比特}
\]

不公平硬币：$p(\text{正}) = 0.9$，$p(\text{反}) = 0.1$
\[
H = -0.9 \log_2 0.9 - 0.1 \log_2 0.1 \approx 0.14 + 0.33 = 0.47 \text{ 比特}
\]

\textbf{结论}：越不公平的硬币，熵越小，结果越“可预测”。
\end{example}

\begin{example}[掷骰子]
公平骰子：$p_i = 1/6$
\[
H = -6 \times \frac{1}{6} \log_2 \frac{1}{6} = \log_2 6 \approx 2.58 \text{ 比特}
\]

作弊骰子（6的概率是50\%，其他各10\%）：
\[
H = -0.5 \log_2 0.5 - 5 \times 0.1 \log_2 0.1 \approx 0.5 + 1.66 = 2.16 \text{ 比特}
\]
\end{example}

\subsection{熵的性质}

\begin{figure}[htbp]
\centering
\begin{tikzpicture}[scale=1.2]
    \begin{axis}[
        xlabel={$p$（正面概率）},
        ylabel={$H(p)$（比特）},
        xmin=0, xmax=1,
        ymin=0, ymax=1.1,
        samples=100,
        domain=0.001:0.999,
        grid=major,
        width=10cm,
        height=6cm,
    ]
    \addplot[very thick, blue] {-x*ln(x)/ln(2) - (1-x)*ln(1-x)/ln(2)};
    
    % 标注最大值点
    \addplot[only marks, mark=*, mark size=3pt, red] coordinates {(0.5, 1)};
    \node[above right, red] at (axis cs:0.5, 0.8) {最大值 $H=1$};
    
    % 标注两端
    \node[below] at (axis cs:0, 0) {\small 确定反面};
    \node[below] at (axis cs:1, 0) {\small 确定正面};
    \end{axis}
\end{tikzpicture}
\caption{二元熵函数 $H(p) = -p\log_2 p - (1-p)\log_2(1-p)$。当 $p=0.5$ 时熵最大，当 $p=0$ 或 $p=1$ 时熵为零。}
\end{figure}

从定义可以直接读出熵的三条基本性质。第一是非负性：$H(p) \geq 0$，等号当且仅当某个 $p_i = 1$，也就是结果完全确定时。第二是最大值：对于 $k$ 个可能结果，$H \leq \log_2 k$，等号当且仅当分布均匀——二元熵函数 $H(p)$ 的曲线在 $p = 0.5$ 处取到最大值 1 比特，正是这条性质最简单的情形。第三是可加性：独立系统的熵等于各部分熵之和，
\[
H(X, Y) = H(X) + H(Y) \quad \text{（若 $X, Y$ 独立）}
\]
这一条直接继承自信息量的可加性。

\begin{insightbox}
同一个量 $H(p)$ 可以从三个角度理解，它们说的是一回事：作为\textbf{不确定性}，熵越大，结果越难预测；作为\textbf{信息量}，熵越大，每次观测带来的平均信息越多；作为\textbf{均匀程度}，熵越大，分布越“平坦”。本章后面的最大熵原理用的主要是第三种理解。
\end{insightbox}

% ------------------------------------------------------------------------------------
\section{熵的组合学起源}

\subsection{微观态计数}

上一节把熵定义成平均信息量，但这个定义看上去仍有几分人为：为什么恰好是 $-\sum p \ln p$，而不是别的某个衡量“平坦程度”的函数？答案来自物理学，而且来自一个纯粹的计数问题。

\textbf{问题}：把 $N$ 个可区分的球放入 $k$ 个盒子，第 $i$ 个盒子有 $n_i$ 个球（$\sum_i n_i = N$）。有多少种放法？

答案是多项式系数：
\begin{equation}
W = \frac{N!}{n_1! n_2! \cdots n_k!}
\end{equation}

\begin{example}[具体计算]
把6个球放入3个盒子，分布为 $(3, 2, 1)$：
\[
W = \frac{6!}{3! \cdot 2! \cdot 1!} = \frac{720}{6 \cdot 2 \cdot 1} = 60 \text{ 种}
\]

如果分布为 $(2, 2, 2)$（均匀）：
\[
W = \frac{6!}{2! \cdot 2! \cdot 2!} = \frac{720}{8} = 90 \text{ 种}
\]

均匀分布有更多的微观实现方式！
\end{example}

\subsection{从计数到熵}

当 $N$ 很大时，用 Stirling 近似 $\ln n! \approx n \ln n - n$：
\clarify{从这里开始，我们改用自然对数 $\ln$（单位是奈特 nat），它与 $\log_2$（比特）只差常数 $\ln 2$。凡是与比特有关的地方（编码、压缩）会再换回 $\log_2$。}

令 $p_i = n_i / N$（各盒子的比例），则：
\begin{align}
\frac{1}{N} \ln W &= \frac{1}{N} \left( \ln N! - \sum_i \ln n_i! \right) \\
&\approx \frac{1}{N} \left( N \ln N - N - \sum_i (n_i \ln n_i - n_i) \right) \\
&= \ln N - \sum_i \frac{n_i}{N} \ln n_i \\
&= -\sum_i p_i \ln p_i
\end{align}

（最后一步用到 $\ln n_i = \ln(p_i N) = \ln p_i + \ln N$）

于是，香农的熵不是凭空定义的：它就是微观态数的对数按球数平摊后的值。这也回答了前面的疑问——$-\sum p \ln p$ 这个形式是从计数里长出来的，不是选出来的。

\begin{keytheorem}[title=熵的组合学意义]
\begin{equation}
H(p) = -\sum_i p_i \ln p_i \approx \frac{1}{N} \ln W
\end{equation}

\textbf{熵 $\propto$ 微观态数的对数}

熵大 $\Leftrightarrow$ 微观实现方式多 $\Leftrightarrow$ 更“自然”、更容易出现
\end{keytheorem}
\storynote{Ludwig Boltzmann}{1844--1906，奥地利物理学家，统计力学奠基人。他的墓碑上刻着$S = k \log W$——熵等于微观态数的对数。这个公式连接了宏观热力学和微观统计力学。他因学术争论而抑郁，最终自杀，但他的理论被后世完全证实。}

\begin{physicalmeaning}
想象把一盒气体分子放到房间里。所有分子挤在角落是低熵状态，只有极少数微观态对应这种宏观状态；分子均匀分布在房间是高熵状态，有海量微观态与之对应。这就是为什么气体会自发扩散——不是因为有什么“力”推动，而是因为均匀分布对应的微观态数压倒性地多！
\end{physicalmeaning}

% ------------------------------------------------------------------------------------
\section{最大熵原理}

\subsection{一个关键问题}

前两节回答了“熵是什么”。现在把问题反过来：能不能用熵去\textbf{构造}分布？假设你在调查某城市的收入分布，但只知道两个信息——平均收入是 5 万元，而且收入必须非负。你应该假设什么样的收入分布？

\begin{figure}[htbp]
\centering
\begin{tikzpicture}[scale=0.9]
    \begin{axis}[
        xlabel={收入（万元）},
        ylabel={概率密度},
        xmin=0, xmax=20,
        ymin=0, ymax=0.25,
        samples=100,
        domain=0:20,
        width=12cm,
        height=6cm,
        legend pos=north east,
    ]
    % 指数分布（最大熵）
    \addplot[very thick, blue] {0.2*exp(-0.2*x)};
    \addlegendentry{指数分布（最大熵）}
    
    % 其他可能的分布
    \addplot[thick, red, dashed] {0.4*x*exp(-0.4*x)};
    \addlegendentry{其他分布1}
    
    \addplot[thick, green!60!black, dotted] {1/(5*sqrt(2*pi))*exp(-(x-5)^2/50)};
    \addlegendentry{其他分布2}
    \end{axis}
\end{tikzpicture}
\caption{满足“均值=5万”约束的不同分布。哪个是“最合理”的选择？}
\end{figure}

有无穷多个分布满足“均值=5万”的条件。选哪一个？

\subsection{最大熵原理的回答}

\begin{quote}
\textit{“不要假装知道你不知道的东西。在约束条件下，选择最均匀（熵最大）的分布。”}
\end{quote}

\begin{keyformula}[title=最大熵原理]
在所有满足已知约束的分布中，选择\textbf{熵最大}的那个。

数学形式：
\begin{equation}
\max_{p} \quad H(p) = -\sum_i p_i \ln p_i
\end{equation}
\begin{align}
\text{s.t.} \quad & \sum_i p_i = 1 \quad \text{（归一化）} \\
& \sum_i p_i g_j(x_i) = c_j, \quad j = 1, \ldots, m \quad \text{（已知约束）}
\end{align}
\end{keyformula}

为什么这是合理的？可以从三个角度回答。从组合学看，上一节的计数表明最大熵分布对应最多的微观态，是“最可能”的宏观状态；从信息论看，最大熵就是最大不确定性，意味着除了已知约束之外不引入任何额外假设；从决策论看，如果你用错误的分布做决策，最大熵分布让你的“最坏情况损失”最小。三条理由指向同一个结论：只知道部分信息时，最大熵分布是最诚实的选择。

\storynote{Edwin Jaynes}{1922--1998，美国物理学家，最大熵原理的现代奠基人。他1957年的论文将熵和概率推理统一起来，影响了统计学、机器学习、贝叶斯推断等领域。他认为``概率是知识的逻辑，不是物理的性质''。}

\begin{tcolorbox}[colback=yellow!10, colframe=orange!80, title=本章的核心观点]
\textbf{最大熵原理不是什么新东西——它就是第1章的拉格朗日乘子法，只不过“变量”从一个向量 $x$ 变成了一个概率分布 $p$，“目标”从某个代价变成了熵。}

约束是什么？——你\textbf{已知}的东西：概率必须归一化，某些量的平均值必须等于观测值。

乘子是什么？——每条已知信息的“熵价格”：$\beta_j = -\partial H^*/\partial c_j$。在物理里它给出温度的倒数，在机器学习里它是特征权重 $w_j$。这不是巧合。

\begin{center}
\renewcommand{\arraystretch}{1.3}
\small
\begin{tabularx}{\linewidth}{llX}
\toprule
 & \textbf{第1章（约束优化）} & \textbf{本章（最大熵）} \\
\midrule
变量 & 向量 $x \in \R^n$ & 分布 $p = (p_1, \ldots, p_k)$ 或密度 $p(x)$ \\
目标 & $\min f(x)$ & $\max H(p) = -\sum_i p_i \ln p_i$ \\
约束 & $g(x) = 0$ & $\sum_i p_i = 1$，$\sum_i p_i g_j(x_i) = c_j$ \\
乘子 & $\lambda$ & $\alpha$（给出配分函数 $Z$），$\beta_j$（特征权重 / 逆温度） \\
一阶条件 & $\nabla f + \lambda \nabla g = 0$ & $-\ln p_i - 1 + \alpha + \sum_j \beta_j g_j(x_i) = 0$ \\
解 & 驻点 $x^*$ & 指数族 $p_i = e^{\sum_j \beta_j g_j(x_i)} / Z$ \\
影子价格 & $\lambda = -df^*/d\epsilon$ & $\beta_j = -\partial H^*/\partial c_j$ \\
\bottomrule
\end{tabularx}
\end{center}
\end{tcolorbox}

\subsection{用拉格朗日乘子法求解}

这是一个带等式约束的优化问题——正是第1章拉格朗日乘子法的舞台！
\review{第1章}{拉格朗日函数 $\Lag = f + \lambda^\top g$，对所有变量（包括乘子）求导令其为零。这里 $f = -H(p)$（最小化负熵），$g$ 是归一化和矩约束，每个约束配一个乘子——本章把它们记作 $\alpha$ 和 $\beta_j$，以免与指数分布的率参数 $\lambda$ 混淆。}

\paragraph{Step 1：构造拉格朗日函数}
\begin{equation}
\mathcal{L}(p, \alpha, \bm{\beta}) = -\sum_i p_i \ln p_i + \alpha\left(\sum_i p_i - 1\right) + \sum_{j=1}^{m} \beta_j \left(\sum_i p_i g_j(x_i) - c_j\right)
\end{equation}

\paragraph{Step 2：对 $p_i$ 求导并令其为零}
\begin{equation}
\frac{\partial \mathcal{L}}{\partial p_i} = -\ln p_i - 1 + \alpha + \sum_j \beta_j g_j(x_i) = 0
\end{equation}

\paragraph{Step 3：解出 $p_i$}
\begin{equation}
\ln p_i = \alpha - 1 + \sum_j \beta_j g_j(x_i)
\end{equation}

\begin{equation}
p_i = \exp\left( \alpha - 1 + \sum_j \beta_j g_j(x_i) \right)
\end{equation}

\paragraph{Step 4：用归一化条件确定 $\alpha$}

由 $\sum_i p_i = 1$，定义配分函数：
\review{第2章}{“$\beta$ 的值由约束隐式确定”其实就是在解对偶问题：$\min_\beta [\ln Z(\beta) - \beta^\top c]$。它无约束、且是凸的（log-sum-exp 减线性），对 $\beta_j$ 求导恰好得到 $\langle g_j \rangle = c_j$。}
\begin{equation}
Z(\bm{\beta}) = \sum_i \exp\left( \sum_j \beta_j g_j(x_i) \right)
\end{equation}

则 $e^{\alpha - 1} = 1/Z$，最终得到：

\begin{keytheorem}[title=最大熵分布的指数形式]
\begin{equation}
\boxed{p_i = \frac{1}{Z} \exp\left( \sum_j \beta_j g_j(x_i) \right)}
\label{eq:maxent}
\end{equation}

其中 $Z = \sum_i \exp\left( \sum_j \beta_j g_j(x_i) \right)$ 是归一化常数（配分函数），$\beta_j$ 是约束 $\langle g_j \rangle = c_j$ 对应的拉格朗日乘子，它的值由约束条件隐式确定。
\end{keytheorem}
\preview{这个“指数族”形式就是深度学习里的 softmax：第8章分类网络的输出层、第11章 Transformer 的注意力权重 $\mathrm{softmax}(q \cdot k / \sqrt{d})$ 都是它的实例，第11章还会把 $\sqrt{d}$ 解释成“温度”。}

\begin{physicalmeaning}
拉格朗日乘子 $\beta_j$ 做的事情是调节分布的形状，使得约束 $\langle g_j \rangle = c_j$ 被精确满足：$\beta_j > 0$ 时，分布倾向于让 $g_j(x)$ 大的状态概率增加；$\beta_j < 0$ 时则相反，倾向于让 $g_j(x)$ 大的状态概率减少。

在物理中，取 $g = E$（能量）时得到玻尔兹曼分布 $p \propto e^{-E/kT}$，即 $\beta = -1/kT$（物理惯例把乘子取相反符号），$|\beta|$ 是逆温度；在机器学习中，$\beta$ 是特征权重。E. T. Jaynes 1957 年提出最大熵原理时，正是把统计力学整个解释成这样一种在已知约束下的推断：玻尔兹曼分布不是额外的物理假设，而是“只知道平均能量”时的最大熵解。

更准确地说：$\beta_j = -\dfrac{\partial H^*}{\partial c_j}$——它度量“多知道一点信息（约束值变一个单位），最大熵下降多少”。这与第1章的 $\lambda = -df^*/d\epsilon$ 完全对应：从“影子价格”到“特征权重”，$\lambda$ 的统一性再次显现！
\end{physicalmeaning}

% ------------------------------------------------------------------------------------
\section{经典分布的推导}

有了式~\eqref{eq:maxent}，推导一个具体的分布就只剩两步：说清楚约束是什么，再由约束把乘子定出来。下面用这套程序推导几个著名的概率分布，约束一次比一次多。
\clarify{如果你没学过概率论，不用担心——本节推导的这些分布（均匀分布、指数分布、高斯分布）是最基础的概率分布，学完本节你就掌握了概率论最核心的几个分布！}

\subsection{均匀分布：一无所知时}

这时唯一的约束是归一化 $\sum_i p_i = 1$，拉格朗日函数只需要一个乘子：
\begin{equation}
\mathcal{L} = -\sum_i p_i \ln p_i + \alpha \left( \sum_i p_i - 1 \right)
\end{equation}

\begin{equation}
\frac{\partial \mathcal{L}}{\partial p_i} = -\ln p_i - 1 + \alpha = 0 \quad \Rightarrow \quad p_i = e^{\alpha - 1}
\end{equation}

由于右边不依赖于 $i$，所以 $p_i$ 对所有 $i$ 相同：
\begin{equation}
\boxed{p_i = \frac{1}{k}} \quad \text{（均匀分布）}
\end{equation}

\begin{insightbox}
当你对一件事\textbf{一无所知}时（没有任何约束），最大熵原理告诉你：\textbf{假设每种可能性相等}。

这就是概率论中的“无差别原理”或“不充分理由原则”。
\end{insightbox}

\subsection{指数分布：已知均值时}

现在多知道一点：变量非负，而且均值已知。约束因此有三条——归一化 $\int_0^\infty p(x) dx = 1$，均值 $\int_0^\infty x \cdot p(x) dx = \mu$，以及支撑 $x \geq 0$（它决定了积分区间）。其中只有一个“矩约束” $g(x) = x$，$c = \mu$。

由最大熵公式：
\begin{equation}
p(x) = \frac{1}{Z} e^{\beta x}
\end{equation}

为了让 $p(x)$ 在 $[0, \infty)$ 上可积，必须 $\beta < 0$。令 $\lambda = -\beta > 0$：
\clarify{这里的 $\lambda = -\beta$ 只是指数分布的率参数（rate），和第1章的乘子 $\lambda$ 不是一回事——本章的乘子记作 $\alpha$、$\beta_j$。}
\begin{equation}
p(x) = \frac{1}{Z} e^{-\lambda x}
\end{equation}

计算归一化常数：
\begin{equation}
Z = \int_0^\infty e^{-\lambda x} dx = \frac{1}{\lambda}
\end{equation}

计算均值约束：
\begin{equation}
\mu = \int_0^\infty x \cdot \lambda e^{-\lambda x} dx = \frac{1}{\lambda}
\end{equation}

所以 $\lambda = 1/\mu$，最终得到：

\begin{equation}
\boxed{p(x) = \frac{1}{\mu} e^{-x/\mu}} \quad \text{（指数分布）}
\end{equation}

\begin{physicalmeaning}
指数分布在放射性衰变（已知平均寿命）、电话呼叫的到达间隔（已知平均间隔）、设备的故障时间（已知平均寿命）里反复出现。这些情况的共同特点是只知道均值，不知道其他信息；最大熵原理告诉我们，这时就应该假设指数分布。
\end{physicalmeaning}
\think{按第1章“影子价格”的定义，$\beta_j$ 应该等于哪个导数？提示：把最优熵 $H^*$ 看成 $c_j$ 的函数，用包络定理，你会得到 $\beta_j = -\partial H^*/\partial c_j$。热力学里 $\partial S/\partial U = 1/T$——温度的倒数正是“能量约束的影子价格”。}

\subsection{高斯分布：已知均值和方差时}

再多知道一点：均值 $\mu$ 和方差 $\sigma^2$ 都已知，支撑是整个实轴。约束是归一化 $\int_{-\infty}^{\infty} p(x) dx = 1$、均值 $\int x \cdot p(x) dx = \mu$ 和方差 $\int (x-\mu)^2 \cdot p(x) dx = \sigma^2$，也就是两个矩约束 $g_1(x) = x$，$g_2(x) = (x-\mu)^2$。

由最大熵公式：
\begin{equation}
p(x) = \frac{1}{Z} \exp\left( \beta_1 x + \beta_2 (x-\mu)^2 \right)
\end{equation}

为了可积，$\beta_2 < 0$。令 $\beta_2 = -1/(2\sigma^2)$，可以验证均值约束要求 $\beta_1 = 0$（因为 $(x-\mu)^2$ 已经以 $\mu$ 为中心）。

最终：
\begin{equation}
\boxed{p(x) = \frac{1}{\sqrt{2\pi\sigma^2}} \exp\left( -\frac{(x-\mu)^2}{2\sigma^2} \right)} \quad \text{（高斯分布）}
\end{equation}
\storynote{高斯}{高斯（1777-1855）被誉为“数学王子”。高斯分布以他命名，但最早是德莫瓦弗在1733年发现的。高斯用它来分析天文观测误差，发现它神奇地描述了几乎所有测量误差的分布。}

\begin{keytheorem}[title=高斯分布的最大熵性质]
在所有均值为 $\mu$、方差为 $\sigma^2$ 的连续分布中，\textbf{高斯分布的熵最大}。

这解释了为什么高斯分布在自然界和统计学中如此普遍——当你只知道均值和方差时，高斯分布是“最诚实”的选择。
\end{keytheorem}

\subsection{经典分布总结}

\begin{center}
\renewcommand{\arraystretch}{1.5}
\begin{tabular}{l|l|l}
\toprule
\textbf{已知约束} & \textbf{最大熵分布} & \textbf{应用场景} \\
\midrule
仅归一化 & 均匀分布 & 完全无知 \\
均值（离散有限） & 几何分布 & 等待时间 \\
均值（连续非负） & 指数分布 & 寿命、间隔 \\
均值 + 方差 & 高斯分布 & 测量误差、自然现象 \\
均值（非负整数） & 泊松分布 & 计数数据 \\
\bottomrule
\end{tabular}
\end{center}
\tryit{用 numpy 各生成 10 万个均值 0、方差 1 的样本（高斯、均匀 $[-\sqrt{3}, \sqrt{3}]$、拉普拉斯），用直方图估计熵，验证高斯最大：理论值分别约 1.42、1.24、1.35 nat。}

% ------------------------------------------------------------------------------------
\section{KL散度：分布之间的“距离”}

\subsection{问题的提出}

到目前为止我们只谈单个分布的熵。但机器学习里更常见的情形是手里有两个分布，需要比较它们的“差异”：模型预测的分布 $q$ 与真实分布 $p$ 有多接近？数据的经验分布与理论分布差多少？一个自然的想法是用“距离”来度量。但什么样的“距离”对概率分布是合理的？

\subsection{KL散度的定义}

\begin{keyformula}[title=KL散度（Kullback-Leibler Divergence）]
\begin{equation}
D_{KL}(P \| Q) = \sum_i p_i \ln \frac{p_i}{q_i} = \sum_i p_i \ln p_i - \sum_i p_i \ln q_i
\end{equation}

或写成期望形式：
\begin{equation}
D_{KL}(P \| Q) = \mathbb{E}_{x \sim P} \left[ \ln \frac{P(x)}{Q(x)} \right]
\end{equation}
\end{keyformula}
\preview{KL散度是深度学习中最重要的概念之一。第8章的神经网络训练、第9章的强化学习策略优化、第12章的大模型后训练（RLHF 里的 KL 约束）、第16章的世界模型，都会反复用到它。现在花时间理解透彻，后面会事半功倍！}

\subsection{KL散度的直觉}

\begin{physicalmeaning}
KL散度的三种理解：

\textbf{1. 编码效率损失}

如果真实分布是 $P$，但你用基于 $Q$ 设计的编码，平均每个符号要多用
\[
D_{KL}(P \| Q) = H(P, Q) - H(P)
\]
比特。这是用“错误模型”的代价。

\textbf{2. 信息增益}

当你从“先验” $Q$ 更新到“后验” $P$，获得的信息量是 $D_{KL}(P \| Q)$。

\textbf{3. 统计区分度}

$D_{KL}(P \| Q)$ 越大，越容易通过采样区分 $P$ 和 $Q$。
\end{physicalmeaning}

\subsection{KL散度的性质}

KL散度最重要的性质是\textbf{非负性}：$D_{KL}(P \| Q) \geq 0$。证明只需要 Jensen 不等式和 $\ln$ 的凹性：
\[
-D_{KL}(P \| Q) = \sum_i p_i \ln \frac{q_i}{p_i} \leq \ln \sum_i p_i \cdot \frac{q_i}{p_i} = \ln \sum_i q_i = \ln 1 = 0
\]
进一步，等号当且仅当 $P = Q$ 时成立，即 $D_{KL}(P \| Q) = 0 \Leftrightarrow P = Q$。这两条使 KL散度可以充当分布之间的“距离”。但它并不是真正的距离——一般 $D_{KL}(P \| Q) \neq D_{KL}(Q \| P)$，交换两个分布的位置，结果会变。

\begin{example}[不对称性的例子]
设 $P = (0.5, 0.5)$，$Q = (0.9, 0.1)$。

\begin{align}
D_{KL}(P \| Q) &= 0.5 \ln \frac{0.5}{0.9} + 0.5 \ln \frac{0.5}{0.1} \\
&= 0.5 \times (-0.588) + 0.5 \times 1.609 = 0.51
\end{align}

\begin{align}
D_{KL}(Q \| P) &= 0.9 \ln \frac{0.9}{0.5} + 0.1 \ln \frac{0.1}{0.5} \\
&= 0.9 \times 0.588 + 0.1 \times (-1.609) = 0.37
\end{align}

$D_{KL}(P \| Q) \neq D_{KL}(Q \| P)$！
\end{example}

\begin{insightbox}
KL散度不对称的直觉：

$D_{KL}(P \| Q)$ 衡量“用 $Q$ 近似 $P$”的代价——在 $P$ 概率大的地方，$Q$ 不能太小。

$D_{KL}(Q \| P)$ 衡量“用 $P$ 近似 $Q$”的代价——在 $Q$ 概率大的地方，$P$ 不能太小。

两个方向关心的“重点区域”不同！
\end{insightbox}

\subsection{最大熵 = 最小KL散度}

\begin{unification}[title=最大熵的等价形式]
\begin{equation}
\max_p H(p) \;\text{s.t. 约束} \quad \Longleftrightarrow \quad \min_p D_{KL}(p \| u) \;\text{s.t. 约束}
\end{equation}

其中 $u$ 是均匀分布。

\textbf{最大熵 = 在约束下离均匀分布最近}
\end{unification}
\preview{把这里的均匀分布 $u$ 换成“旧策略”，$\min D_{KL}(p \| u)$ 就变成第9章信任域方法（TRPO/PPO）里“新策略别离旧策略太远”的约束。}

\textit{证明}：
\begin{align}
D_{KL}(p \| u) &= \sum_i p_i \ln \frac{p_i}{1/k} = \sum_i p_i \ln p_i + \sum_i p_i \ln k \\
&= -H(p) + \ln k
\end{align}

由于 $\ln k$ 是常数，$\min D_{KL}(p \| u) \Leftrightarrow \max H(p)$。

\section{KL散度的第一性原理}

\subsection{什么是“好的”分布距离？}

假设我们想用一个简单的分布 $q(x)$ 去近似复杂的真实分布 $p(x)$。我们需要一个“差异度量” $D(p \| q)$ 来衡量近似的好坏。

但是，什么样的度量才是“合理”的？信息论给出了答案：合理的度量必须满足三条公理。

\begin{keyformula}[title={信息距离的三条公理 (Shore \& Johnson, 1980)}]
\begin{enumerate}
    \item \textbf{局部性 (Locality)}：差异应该是逐点累积的，不依赖全局耦合。
    \[
    D(p \| q) = \int f(p(x), q(x)) \, dx
    \]
    
    \item \textbf{可加性 (Additivity)}：对于独立系统 $A, B$，联合分布的差异等于各自差异之和。
    \[
    D(p_A p_B \| q_A q_B) = D(p_A \| q_A) + D(p_B \| q_B)
    \]
    
    \item \textbf{一致性 (Consistency)}：当 $q(x) = p(x)$ 时，差异为零且达到最小值。
\end{enumerate}
\end{keyformula}

\begin{physicalmeaning}
这三条公理各有直觉。局部性说的是每个点的偏差独立计算，就像质量、能量一样是“广延量”；可加性说的是独立系统的总偏差等于各自偏差之和，这与熵、能量的性质一致；一致性说的是完美近似时差异为零——这是最基本的要求。
\end{physicalmeaning}

\subsection{从公理推导KL散度}

\paragraph{Step 1：利用可加性锁定函数形式}

设差异度量为 $D(p\|q) = \sum_i h(p_i, q_i)$。

考虑两个独立事件，联合概率 $p_{ij} = u_i v_j$，$q_{ij} = \alpha_i \beta_j$。根据可加性：
\[
\sum_{i,j} h(u_i v_j, \alpha_i \beta_j) = \sum_i h(u_i, \alpha_i) + \sum_j h(v_j, \beta_j)
\]

这是一个经典的\textbf{柯西函数方程}变体。数学上可以证明，唯一满足此性质的连续函数形式必须包含\textbf{对数项}：
\[
h(p, q) = c \cdot p \ln \frac{p}{q} + \cdots
\]

\paragraph{Step 2：利用一致性确定系数}

设通解为：
\[
D(p\|q) = \int \left( A p(x) \ln p(x) + B p(x) \ln q(x) + C p(x) \right) dx
\]

根据一致性公理：当 $q(x) \to p(x)$ 时，$D$ 必须达到最小值（且为0）。

对 $q(x)$ 求变分导数（考虑归一化约束）：
\[
\frac{\delta}{\delta q} \left( D(p\|q) + \lambda \int q \, dx \right) = 0
\quad \Rightarrow \quad
\frac{B p(x)}{q(x)} + \lambda = 0
\]

为了让最优解是 $q^*(x) = p(x)$、且 $D(p\|p) = 0$，需要 $A + B = 0$，$B < 0$（极小性）；取归一化 $A = -B = 1$（相当于选定对数底），就得到：

\begin{keytheorem}[title=KL散度的唯一性]
满足局部性、可加性、一致性三条公理的差异度量，在相差一个正常数（即对数底的选择）的意义下是\textbf{唯一}的：
\begin{equation}
\boxed{D_{KL}(p \| q) = \int p(x) \ln \frac{p(x)}{q(x)} \, dx = -H(p) - \int p(x) \ln q(x) \, dx}
\end{equation}
\end{keytheorem}

\begin{insightbox}
\textbf{为什么机器学习中常用交叉熵损失？}

最小化 $D_{KL}(p^* \| q)$ 等价于：
\[
\min_q \left( \underbrace{-H(p^*)}_{\text{常数，与 }q\text{ 无关}} + \underbrace{\left(-\int p^*(x) \ln q(x) \, dx\right)}_{\text{交叉熵 }H(p^*, q)} \right)
\]

因此，\textbf{最小化交叉熵 = 最小化KL散度 = 找最逼近真实分布的模型}。

这不是“选择”，而是\textbf{数学必然}！
\end{insightbox}

% ------------------------------------------------------------------------------------
\section{常见损失函数的“违规记录”}

既然 KL散度是唯一满足三条公理的度量，为什么实践中还有那么多别的损失函数？因为并非所有损失函数都需要满足信息论公理。有些为了特定的工程目的（如鲁棒性、几何感知），故意违反某些公理。

\begin{center}
\renewcommand{\arraystretch}{1.6}
\begin{tabularx}{\textwidth}{p{2.4cm}|p{2.8cm}|p{2cm}|X}
\toprule
\textbf{损失函数} & \textbf{数学形式} & \textbf{违反的公理} & \textbf{为何要违反？} \\
\midrule
\textbf{交叉熵} (Cross Entropy) 
& $-\sum p \ln q$ 
& 无 
& 信息论的“正统”选择 \\
\midrule
\textbf{MSE} (均方误差) 
& $\int (p-q)^2 dx$ 
& 可加性 
& 计算便利，求导简单，适合回归任务 \\
\midrule
\textbf{MAE} (L1 Loss) 
& $\int |p-q| dx$ 
& 可加性 
& 对异常值鲁棒，不被离群点拉偏 \\
\midrule
\textbf{Hinge} (SVM) 
& $\max(0, 1-yf(x))$ 
& 一致性 
& 追求稀疏性和分类边界，不追求估计真实概率 \\
\midrule
\textbf{Wasserstein} (WGAN) 
& $\inf \mathbb{E}[\|x-y\|]$ 
& 局部性 
& 引入空间距离（全局耦合），解决分布不重叠时KL散度梯度消失的问题 \\
\bottomrule
\end{tabularx}
\end{center}

\begin{physicalmeaning}
\textbf{选择损失函数的智慧}：如果你相信数据服从某个概率模型，就用交叉熵，它是信息论意义下的最优选择；如果你想要对异常值鲁棒，就用 MAE 或 Huber 损失；如果你只关心分类边界，就用 Hinge 损失；如果两个分布可能不重叠，就用 Wasserstein 距离。没有“最好”的损失函数，只有最适合问题的损失函数——每一次“违规”都是在用信息论最优性换取某种工程性质。
\end{physicalmeaning}


% ------------------------------------------------------------------------------------
\section{交叉熵与机器学习}
\preview{本节会出现“神经网络”、“分类器”等深度学习术语。如果你没学过深度学习，暂时只需知道：神经网络是一种能从数据中学习规律的数学模型，分类器是判断“这张图是猫还是狗”之类问题的程序。第8章会详细介绍。}

\subsection{交叉熵的定义}

\begin{keyformula}[title=交叉熵]
\begin{equation}
H(P, Q) = -\sum_i p_i \ln q_i = H(P) + D_{KL}(P \| Q)
\end{equation}

交叉熵 = 真实分布的熵 + KL散度
\end{keyformula}

\subsection{为什么机器学习用交叉熵损失？}

前面的推导都是关于分布本身的，现在把它接到机器学习最常见的任务——分类——上。在分类问题中，$P$ 是真实标签分布（通常是 one-hot，如 $[0, 1, 0]$），$Q$ 是模型预测的概率分布（如 $[0.1, 0.7, 0.2]$）。训练目标是让 $Q$ 尽量接近 $P$，即最小化 $D_{KL}(P \| Q)$。

由于 $H(P)$ 是固定的（真实标签不变），所以：
\begin{equation}
\min_Q D_{KL}(P \| Q) \quad \Longleftrightarrow \quad \min_Q H(P, Q)
\end{equation}

对于 one-hot 标签 $y = [0, \ldots, 1, \ldots, 0]$（第 $c$ 类为1）：
\begin{equation}
H(P, Q) = -\sum_i y_i \ln q_i = -\ln q_c
\end{equation}

这就是我们熟悉的\textbf{交叉熵损失}！

\begin{insightbox}
\textbf{为什么不用均方误差（MSE）？}对于分类问题，交叉熵有三个优势。它是信息论最优的：交叉熵直接度量“用模型编码真实数据”的代价。它的梯度性质好：当预测很错时，交叉熵的梯度大，而 MSE 的梯度可能很小（梯度消失）。它还有概率解释：最小化交叉熵等价于最大化似然。
\end{insightbox}

\subsection{Softmax = 最大熵分布}

神经网络分类器的输出通常经过 softmax 函数：
\begin{equation}
p_i = \frac{e^{z_i}}{\sum_j e^{z_j}}
\end{equation}
\clarify{softmax函数把一组任意实数变成概率分布（全为正数且和为1）。它是深度学习中最常用的函数之一。在第8章你会看到它如何用于图像分类、第11章会看到它如何用于Transformer注意力机制。}

这个形式是不是很眼熟？让我们一步步看它是如何从最大熵原理推导出来的。

\begin{insightbox}
\textbf{Softmax的推导：从最大熵到神经网络}

回顾前面推导的最大熵分布形式（式~\eqref{eq:maxent}）：
\[
p_i = \frac{1}{Z} \exp\left( \sum_j \beta_j g_j(x_i) \right)
\]

现在考虑分类问题：我们想预测输入 $\bm{x}$ 属于哪个类别。

\textbf{Step 1}：神经网络为每个类别 $i$ 计算一个得分（logit）：
\[
z_i = \bm{w}_i^\top \bm{x} + b_i
\]

\textbf{Step 2}：将这个得分代入最大熵分布。这里 $z_i$ 就是 $\sum_j \beta_j g_j(x_i)$ 的具体形式——它表示类别 $i$ 的“能量”或“偏好程度”。

\textbf{Step 3}：归一化得到概率：
\[
p_i = \frac{e^{z_i}}{\underbrace{\sum_j e^{z_j}}_{Z = \text{配分函数}}}
\]

这正是 softmax 函数！
\end{insightbox}

\begin{unification}[title=Softmax的最大熵解释]
Softmax 函数给出的是满足约束 $\langle z \rangle = \sum_i p_i z_i = c$ 的\textbf{最大熵分布}。

其中 $z_i$ 是神经网络的 logits（未归一化得分），$c$ 是由网络参数决定的常数。

\textbf{神经网络学习的权重 $\bm{w}_i$ 是最大熵分布的拉格朗日乘子，logit $z_i = \bm{w}_i^\top \bm{x}$ 是“乘子 $\times$ 特征”。}
\end{unification}

\begin{physicalmeaning}
\textbf{为什么叫 “softmax”？}

考虑一个“硬”版本——argmax，它只输出得分最高的类别：
\[
\text{argmax: } p_i = \begin{cases} 1 & \text{if } i = \arg\max_j z_j \\ 0 & \text{otherwise} \end{cases}
\]

Softmax 是它的“软”版本：当各 $z_i$ 差异很大时，softmax 接近 argmax（赢者通吃）；当各 $z_i$ 差异很小时，softmax 接近均匀分布（不确定）。

你可以通过引入“温度”参数 $T$ 来控制软硬程度：
\[
p_i = \frac{e^{z_i / T}}{\sum_j e^{z_j / T}}
\]
$T \to 0$：变“硬”（接近argmax）；$T \to \infty$：变“软”（接近均匀分布）。

这与物理中的温度概念完全对应！高温时系统混乱（高熵），低温时系统有序（低熵）。
\end{physicalmeaning}

\subsection{逻辑回归 = 最大熵分类}

逻辑回归的模型：
\begin{equation}
P(y=1 | \bm{x}) = \sigma(\bm{w}^\top \bm{x}) = \frac{1}{1 + e^{-\bm{w}^\top \bm{x}}} = \frac{e^{\bm{w}^\top \bm{x}}}{1 + e^{\bm{w}^\top \bm{x}}}
\end{equation}

这可以理解为：在特征期望约束 $\mathbb{E}_{model}[\bm{\phi}(\bm{x})] = \mathbb{E}_{data}[\bm{\phi}(\bm{x})]$ 下的最大熵分布。

权重 $\bm{w}$ 正是这些约束的拉格朗日乘子！

\begin{physicalmeaning}
逻辑回归的训练过程可以完全用最大熵的语言复述：先从数据计算特征的经验期望 $\mathbb{E}_{data}[\bm{\phi}]$，然后调整 $\bm{w}$ 使得模型的特征期望 $\mathbb{E}_{model}[\bm{\phi}]$ 与之匹配，而在满足这些矩约束的同时，模型保持熵最大，不做额外假设。这就是为什么逻辑回归又叫“最大熵分类器”！
\end{physicalmeaning}

% ------------------------------------------------------------------------------------




\section{信息压缩：熵的工程应用}

\subsection{为什么需要变长编码？}

熵被定义成“平均信息量”，这个说法能不能落到实处——它到底对应多少个比特？信源编码给出了最直接的回答。假设你要传输天气预报，有4种天气：晴、阴、雨、雪。

\textbf{方案1：等长编码}

每种天气用2位二进制表示：
\begin{center}
\begin{tabular}{c|c}
天气 & 编码 \\
\hline
晴 & 00 \\
阴 & 01 \\
雨 & 10 \\
雪 & 11 \\
\end{tabular}
\end{center}

传输100天的天气，需要 $100 \times 2 = 200$ 比特。

\textbf{但是}，如果这个城市70\%的日子是晴天呢？

用等长编码太浪费了！“晴”出现70次，每次都用2比特；而“雪”可能只出现5次。

\begin{insightbox}
\textbf{核心直觉}：高频符号用短码，低频符号用长码。

就像常用汉字笔画少（“一、人、大”），罕见汉字笔画多（“{\rarecjkfont 龘、靐}、鑫”）。

或者想想摩尔斯电码：最常用的字母 E 只需要一个点（$\cdot$），而罕见的 Q 需要四个符号（$--\cdot-$）。
\end{insightbox}

\subsection{变长编码的陷阱：前缀码}

变长编码有个问题：\textbf{怎么知道一个符号在哪里结束？}

\begin{example}[歧义编码]
假设编码是：A = 0, B = 01, C = 1

收到 “01”，是 “AB” 还是 “B”？无法确定！
\end{example}

解决方案是\textbf{前缀码}：任何符号的编码都不是另一个符号编码的前缀。

\begin{example}[前缀码]
A = 0, B = 10, C = 11。收到 “01011” 时，第一位是 0，只能是 A（没有其他编码以 0 开头然后继续）；接下来的 10 是 B；最后的 11 是 C。解码结果 ABC，没有歧义！
\end{example}

\begin{physicalmeaning}
前缀码可以用\textbf{二叉树}表示：每个符号对应一个叶子节点，从根到叶子的路径就是编码（左=0，右=1）。因为符号都在叶子上，任何编码都不会是另一个编码的前缀，所以不会有歧义。
\end{physicalmeaning}

\subsection{最优编码的理论极限}

那么，最好的变长编码能做到多好？答案正是熵。

\begin{keytheorem}[title=Shannon第一定理（信源编码定理）]
对于信源 $X$，任何无损前缀码的平均码长都不低于 $H(X)$，而\textbf{最优}前缀码的平均码长 $\bar{L}$ 满足：
\begin{equation}
H(X) \leq \bar{L} < H(X) + 1
\end{equation}

\textbf{熵是压缩的理论极限}——不可能把平均码长压到熵以下！
\end{keytheorem}
\think{为什么平均码长不可能低于熵？提示：把码长看成一个分布 $q_i = 2^{-l_i}$（Kraft 不等式保证 $\sum q_i \le 1$），平均码长 $\sum p_i l_i = -\sum p_i \log_2 q_i$ 就是交叉熵 $H(p, q) \ge H(p)$——正是习题 5 的 Gibbs 不等式。}

直觉上，如果一个符号的概率是 $p_i$，它携带的信息量是 $-\log_2 p_i$ 比特，编码这个符号就至少需要 $-\log_2 p_i$ 比特，否则就“凭空创造”了信息；平均下来，每个符号至少需要 $H = \sum p_i (-\log_2 p_i)$ 比特。理想情况下，符号 $i$ 的最优码长应该是：
\begin{equation}
l_i^* = -\log_2 p_i
\end{equation}

但码长必须是整数！所以实际码长是 $\lceil -\log_2 p_i \rceil$，这就是为什么有 “+1” 的间隙。

\subsection{Huffman编码：贪心构造最优码}

现在的问题是：如何构造一个平均码长最小的前缀码？

\textbf{关键洞察}：前缀码对应二叉树，平均码长就是叶子的加权平均深度：
\begin{equation}
\bar{L} = \sum_i p_i \cdot \text{depth}(i)
\end{equation}

我们要最小化这个加权深度！

\begin{physicalmeaning}
想象一下：你要把符号放在一棵二叉树的叶子上。高频符号应该放在\textbf{靠近根}的位置（深度小，编码短），低频符号可以放在\textbf{远离根}的位置（深度大，编码长）。就像公司组织架构：CEO（最重要）在顶层，实习生在最底层。
\end{physicalmeaning}

\textbf{Huffman的贪心策略}：从叶子往上建树，每次把概率最小的两个节点合并成一个。1951 年，David Huffman 还是 MIT 的研究生，导师 Fano 允许学生用一篇论文代替期末考试，题目正是最优二进制编码；他在几乎放弃的时候想到了这个“自底向上合并最小概率”的办法。

为什么这样做是对的？直觉是：
\begin{quote}
概率最小的两个符号，应该在树的最深处，而且应该是“兄弟”（共享同一个父节点）。因为它们最不重要，放在最深处损失最小。
\end{quote}

\subsection{Huffman编码的详细例子}

\begin{example}[天气预报编码]
某城市的天气概率：
\begin{center}
\begin{tabular}{c|cccc}
天气 & 晴 & 阴 & 雨 & 雪 \\
\hline
概率 & 0.5 & 0.25 & 0.15 & 0.1 \\
\end{tabular}
\end{center}

\textbf{Step 1}：找概率最小的两个——雨(0.15)和雪(0.1)，合并
\begin{center}
\begin{tikzpicture}[scale=0.8]
    \node[draw, circle, minimum size=8mm] (rs) at (0,0) {0.25};
    \node[draw, circle, fill=blue!15, minimum size=8mm] (r) at (-1,-1.2) {雨};
    \node[draw, circle, fill=blue!15, minimum size=8mm] (s) at (1,-1.2) {雪};
    \draw (rs) -- node[left, font=\scriptsize] {0} (r);
    \draw (rs) -- node[right, font=\scriptsize] {1} (s);
    \node[below, font=\scriptsize, gray] at (r.south) {0.15};
    \node[below, font=\scriptsize, gray] at (s.south) {0.1};
\end{tikzpicture}
\end{center}

现在有：晴(0.5)、阴(0.25)、[雨雪](0.25)

\textbf{Step 2}：找最小的两个——阴(0.25)和[雨雪](0.25)，合并
\begin{center}
\begin{tikzpicture}[scale=0.8]
    \node[draw, circle, minimum size=8mm] (top) at (0,0) {0.5};
    \node[draw, circle, fill=blue!15, minimum size=8mm] (y) at (-1.5,-1.2) {阴};
    \node[draw, circle, minimum size=8mm] (rs) at (1.5,-1.2) {0.25};
    \node[draw, circle, fill=blue!15, minimum size=8mm] (r) at (0.5,-2.4) {雨};
    \node[draw, circle, fill=blue!15, minimum size=8mm] (s) at (2.5,-2.4) {雪};
    \draw (top) -- node[left, font=\scriptsize] {0} (y);
    \draw (top) -- node[right, font=\scriptsize] {1} (rs);
    \draw (rs) -- node[left, font=\scriptsize] {0} (r);
    \draw (rs) -- node[right, font=\scriptsize] {1} (s);
    \node[below, font=\scriptsize, gray] at (y.south) {0.25};
    \node[below, font=\scriptsize, gray] at (r.south) {0.15};
    \node[below, font=\scriptsize, gray] at (s.south) {0.1};
\end{tikzpicture}
\end{center}

现在有：晴(0.5)、[阴雨雪](0.5)

\textbf{Step 3}：合并最后两个——晴(0.5)和[阴雨雪](0.5)
\begin{center}
\begin{tikzpicture}[scale=0.8]
    \node[draw, circle, minimum size=8mm] (root) at (0,0) {1.0};
    \node[draw, circle, fill=blue!15, minimum size=8mm] (q) at (-3,-1.2) {晴};
    \node[draw, circle, minimum size=8mm] (right) at (2,-1.2) {0.5};
    \node[draw, circle, fill=blue!15, minimum size=8mm] (y) at (0.5,-2.4) {阴};
    \node[draw, circle, minimum size=8mm] (rs) at (3.5,-2.4) {0.25};
    \node[draw, circle, fill=blue!15, minimum size=8mm] (r) at (2.5,-3.6) {雨};
    \node[draw, circle, fill=blue!15, minimum size=8mm] (s) at (4.5,-3.6) {雪};
    
    \draw (root) -- node[left, font=\scriptsize] {0} (q);
    \draw (root) -- node[right, font=\scriptsize] {1} (right);
    \draw (right) -- node[left, font=\scriptsize] {0} (y);
    \draw (right) -- node[right, font=\scriptsize] {1} (rs);
    \draw (rs) -- node[left, font=\scriptsize] {0} (r);
    \draw (rs) -- node[right, font=\scriptsize] {1} (s);
    
    \node[below, font=\scriptsize, gray] at (q.south) {0.5};
    \node[below, font=\scriptsize, gray] at (y.south) {0.25};
    \node[below, font=\scriptsize, gray] at (r.south) {0.15};
    \node[below, font=\scriptsize, gray] at (s.south) {0.1};
    
    % 编码标注
    \node[left=5pt, font=\footnotesize, blue] at (q.west) {编码: 0};
    \node[left=5pt, font=\footnotesize, blue] at (y.west) {编码: 10};
    \node[below left, font=\footnotesize, blue] at (r.south west) {编码: 110};
    \node[below right, font=\footnotesize, blue] at (s.south east) {编码: 111};
\end{tikzpicture}
\end{center}

\textbf{最终编码}：
\begin{center}
\begin{tabular}{c|c|c|c|c}
天气 & 概率 & 编码 & 码长 & 理想码长 $-\log_2 p$ \\
\hline
晴 & 0.5 & 0 & 1 & 1.0 \\
阴 & 0.25 & 10 & 2 & 2.0 \\
雨 & 0.15 & 110 & 3 & 2.74 \\
雪 & 0.1 & 111 & 3 & 3.32 \\
\end{tabular}
\end{center}

\textbf{性能分析}：平均码长为
\[
\bar{L} = 0.5 \times 1 + 0.25 \times 2 + 0.15 \times 3 + 0.1 \times 3 = 1.75
\]
比特，而熵为
\[
H = -0.5\log_2 0.5 - 0.25\log_2 0.25 - 0.15\log_2 0.15 - 0.1\log_2 0.1 \approx 1.74
\]
比特，效率 $\eta = H / \bar{L} = 1.74 / 1.75 = 99.4\%$。

对比等长编码（2比特/符号），Huffman编码节省了 $(2-1.75)/2 = 12.5\%$！

传输100天天气：等长需要200比特，Huffman只需要175比特。
\end{example}

\begin{insightbox}
\textbf{Huffman编码为什么是最优的？}直觉证明分三步：概率最小的两个符号必须在最深层，否则可以把它们交换到更深处、减小平均码长；它们必须是兄弟，否则可以把它们配对、腾出其他位置；合并之后问题规模减一，对剩下的问题递归应用同样的论证。这是一个经典的\textbf{贪心算法}：每步做局部最优选择（合并最小两个），最终得到全局最优。
\end{insightbox}

\subsection{从Huffman到现代压缩}

Huffman编码是所有现代压缩算法的基础：

\begin{center}
\renewcommand{\arraystretch}{1.4}
\begin{tabular}{l|l|l}
\toprule
\textbf{算法} & \textbf{核心思想} & \textbf{应用} \\
\midrule
Huffman & 静态概率，构建最优树 & JPEG, MP3 的一部分 \\
算术编码 & 把整个消息编码为一个小数 & 更接近熵极限 \\
LZ77/LZ78 & 利用重复模式（字典压缩） & ZIP, GZIP, PNG \\
LZW & LZ78的变种 & GIF, Unix compress \\
Deflate & LZ77 + Huffman & ZIP, GZIP, HTTP \\
\bottomrule
\end{tabular}
\end{center}

\begin{physicalmeaning}
\textbf{为什么ZIP能把文件压小？}因为大多数文件都有\textbf{冗余}：文本中某些字母和词组频繁出现（字母 e 比 z 常见100倍），代码中有大量重复模式（括号、关键字），图像中相邻像素通常很相似。压缩算法就是在\textbf{发现并利用这些冗余}，用更少的比特表示同样的信息。

理论极限就是熵——当且仅当数据完全随机（无冗余）时，才无法压缩。
\end{physicalmeaning}

\subsection{实际文本的压缩潜力}

\begin{center}
\renewcommand{\arraystretch}{1.3}
\begin{tabular}{l|c|l}
\toprule
\textbf{模型} & \textbf{熵率（bits/字母）} & \textbf{说明} \\
\midrule
完全随机（26字母均匀） & $\log_2 26 \approx 4.7$ & 无任何结构 \\
单字母频率 & $\approx 4.2$ & 考虑 e 比 z 常见 \\
二元组（bigram） & $\approx 3.5$ & 考虑字母对频率 \\
人类实验估计 & $\approx 1.0 - 1.5$ & 利用全部语言知识 \\
GPT-4 估计 & $\approx 0.9$ & 当前最好的语言模型 \\
\bottomrule
\end{tabular}
\end{center}

\begin{insightbox}
英语文本的真实熵率约为 1-1.5 bits/字母，远低于 4.7 bits 的“随机上限”。

这意味着\textbf{英语有约 70\% 的冗余}！正是这些冗余，让文本压缩（ZIP, GZIP）能够工作，让语言模型（GPT）能预测下一个词，让手机自动补全能猜对，也让你能读懂 “Cn y rd ths sntnc?”。
\end{insightbox}
% ------------------------------------------------------------------------------------
\section{本章小结}

\begin{enumerate}
    \item \textbf{信息量}：事件 $x$ 的信息量是 $I(x) = -\log_2 p(x)$，概率越小，信息量越大。

    \item \textbf{熵}：$H(p) = -\sum p_i \log p_i$ 是平均信息量，度量一个分布的不确定性。

    \item \textbf{最大熵原理}：在约束下选择熵最大的分布，用拉格朗日乘子法解出来就是指数族形式
    \[
    p_i = \frac{1}{Z} \exp\left( \sum_j \beta_j g_j(x_i) \right) .
    \]

    \item \textbf{拉格朗日乘子}：$\beta_j$ 是约束的拉格朗日乘子，它调节分布的形状，使约束被精确满足。

    \item \textbf{KL散度}：$D_{KL}(P \| Q) = \sum p_i \ln \frac{p_i}{q_i}$ 度量两个分布的差异，非负但不对称。

    \item \textbf{交叉熵}：$H(P, Q) = H(P) + D_{KL}(P \| Q)$ 是机器学习的核心损失函数。

    \item \textbf{机器学习联系}：softmax 就是最大熵分布，逻辑回归就是最大熵分类，而最小化交叉熵、最小化KL散度和最大似然是同一件事。

    \item \textbf{信息压缩}：熵是压缩的理论极限，最优码长约为 $-\log_2 p$。
\end{enumerate}

\begin{unification}[title=本章的统一视角]
\begin{center}
\begin{tabular}{c|c}
\toprule
\textbf{信息论概念} & \textbf{优化视角} \\
\midrule
熵 $H(p)$ & 目标函数（最大化） \\
归一化 $\sum p_i = 1$ & 等式约束 \\
矩约束 $\langle g \rangle = c$ & 等式约束 \\
拉格朗日乘子 $\beta$ & 特征权重 / 逆温度 \\
最大熵分布 & 拉格朗日条件的解 \\
KL散度 & 与均匀分布的“距离” \\
\bottomrule
\end{tabular}
\end{center}

信息论的核心问题——“在约束下找最均匀的分布”——正是第1章拉格朗日乘子法的直接应用！
\end{unification}

\begin{physicalmeaning}
\textbf{$\lambda$ 在本章的意义}：本章的乘子叫 $\beta_j$，它是第 $j$ 条已知信息的“熵价格”——$\beta_j = -\partial H^*/\partial c_j$，约束值每变一个单位，最大可达的熵变多少。在物理里，$\partial S/\partial U = 1/T$，温度的倒数是“能量约束的影子价格”；在机器学习里，logit $z = w^\top x$ 中的权重 $w$ 就是乘子，softmax 就是最大熵分布。往前看，第1章的影子价格、第5章的约束力、本章的逆温度，是同一个 $\lambda$；第8章会把 softmax 放进网络的最后一层，第11章会把它变成注意力权重，第9章的最大熵强化学习（SAC）则把熵直接加进目标。
\end{physicalmeaning}

% ------------------------------------------------------------------------------------
\section{习题}

\paragraph{计算题}
\begin{enumerate}
    \item 计算以下分布的熵：
    \begin{enumerate}
        \item 公平硬币：$p = (0.5, 0.5)$
        \item 偏硬币：$p = (0.8, 0.2)$
        \item 公平骰子：$p = (1/6, 1/6, 1/6, 1/6, 1/6, 1/6)$
    \end{enumerate}
    
    \item 计算 $D_{KL}(P \| Q)$ 和 $D_{KL}(Q \| P)$，其中 $P = (0.5, 0.5)$，$Q = (0.8, 0.2)$。验证它们不相等。
    
    \item 为符号集 $\{A, B, C, D\}$（概率分别为 $0.5, 0.25, 0.125, 0.125$）构造 Huffman 编码，计算平均码长，并与熵比较。
\end{enumerate}

\paragraph{证明题}
\begin{enumerate}
    \setcounter{enumi}{3}
    \item 证明：在只有归一化约束下，最大熵分布是均匀分布。
    
    \item 证明 Gibbs 不等式：$D_{KL}(P \| Q) \geq 0$，等号成立当且仅当 $P = Q$。
    
    （提示：使用 Jensen 不等式和 $\ln x \leq x - 1$）
    
    \item 证明：对于固定方差 $\sigma^2$，高斯分布是连续分布中熵最大的。
    
    （提示：设任意分布 $p(x)$，计算 $H(p) - H(\text{Gaussian})$ 并说明它 $\leq 0$）
\end{enumerate}

\paragraph{编程题}
\begin{enumerate}
    \setcounter{enumi}{6}
    \item 实现 Huffman 编码算法：
    \begin{enumerate}
        \item 输入：符号概率分布
        \item 输出：每个符号的编码
        \item 验证：平均码长接近熵
    \end{enumerate}
    
    \item 实现最大熵分类器（softmax 回归）：
    \begin{enumerate}
        \item 使用 MNIST 数据集
        \item 实现前向传播（softmax）和反向传播
        \item 使用交叉熵损失
        \item 报告分类准确率
    \end{enumerate}
\end{enumerate}

\paragraph{思考题}
\begin{enumerate}
    \setcounter{enumi}{8}
    \item 为什么深度学习中分类任务用交叉熵损失而不是均方误差？从梯度的角度分析。
    
    \item KL散度不对称意味着什么？在实际应用中，对 $Q$ 最小化 $D_{KL}(P \| Q)$ 与最小化 $D_{KL}(Q \| P)$，分别倾向于产生什么样的近似？
    
    \item 最大熵原理与 Occam 剃刀有什么联系？为什么“最均匀”的分布是“最简单”的？
    
    \item GPT 等语言模型本质上在估计什么？它与熵率有什么关系？
\end{enumerate}